\documentclass[10pt,letterpaper]{article}
\usepackage{cornell-notes}

\title{Clause 27.11.06: Uninitialized Move}
\author{}
\date{}
\setDocTitle{Clause 27.11.06: Uninitialized Move}
\setDocSubtitle{C++ 2024}
\setDocOwner{C++ 2024 Cornell Notes}
\setDocDate{\today}
\setCornellCollection{C++ 2024 Cornell Notes}
\setCornellUnitType{Clause}
\setCornellUnitNumber{27.11.06}
\setCornellUnitTitle{Uninitialized Move}
\hypersetup{pdftitle={Clause 27.11.06: Uninitialized Move},pdfsubject={Cornell notes},pdfauthor={}}

\begin{document}
\maketitle
\begin{CornellOverviewBox}{Overview}
\textbf{Classification:} Standard library - Normative\\[2pt]
\textbf{Learning objective:} Understand the rules, terminology, and semantic consequences associated with uninitialized\_move.
\end{CornellOverviewBox}\begin{CornellSummaryBox}{Summary}
{\sffamily\bfseries\color{Accent} Summary}\par\vspace{3pt}Section 27.11.6 focuses on uninitialized\_move. Apply the specified interface together with its mandates, effects, result conditions, exception behavior, and complexity constraints. When applying it, preserve the distinction between mandatory requirements, recommendations, permissions, and behavior deliberately left undefined, unspecified, or implementation-defined.
\end{CornellSummaryBox}

\end{document}
